% Slide 3 – what the prototype already does (criterion: quality and completeness of the prototype, 20 %).
\begin{frame}{A working prototype: a web app for phone and desktop}
  \begin{columns}[T,onlytextwidth]
    \begin{column}{0.57\textwidth}
      \small
      \begin{itemize}
        \setlength{\itemsep}{2.5mm}
        \item \textbf{6,545 places, 27 accessibility features}\\
              Concrete information about the entrance, lift, toilet, parking and other needs.
        \item \textbf{Matched to the user}\\
              The needs profile filters and ranks places by what really matters.
        \item \textbf{Live, trustworthy data}\\
              Source, date and verification status for every piece of information, plus community confirmations.
        \item \textbf{AI + barrier map}\\
              The assistant recommends places, and users report and update real obstacles.
      \end{itemize}
    \end{column}
    \begin{column}{0.41\textwidth}
      \screenshot{mapa-desktop}
      \vspace{0.3mm}
      \muted{“Wheelchair” profile: places, layers, reported barriers with severity and confirmations.
        The list next to the map is its text alternative.}
      \vspace{1.5mm}

      \kpi{6,545}{places from OSM}\hfill\kpi{11}{data layers}\hfill\kpi{503}{tests, offline}
    \end{column}
  \end{columns}

  \note[item]{This is not a mock-up: everything on the slide runs from the repository (Python standard library plus plain JavaScript, no dependencies).}
  \note[item]{Place data from OSM: 1,495 of 6,545 places already have at least one accessibility attribute; the rest show “No data”, and that is exactly what we display.}
  \note[item]{Layers: stops with platform kerb type, toilets, disabled parking, lifts, kerbs, crossings with tactile paving, bike routes and infrastructure, luggage storage, A\&E and 24-hour pharmacies.}
  \note[item]{The AI assistant is optional; the demo container runs the keyword mode, so nothing depends on the network.}
\end{frame}
